% ============================================================
% 数学\统计图表\六年级.tex
% 本文件由 main.tex 通过 \input 引入，请勿单独编译
% ============================================================


\section{解答：六年级一班女生1分钟跳绳成绩统计表}

\vspace{0.4cm}

\noindent\textbf{1. 题目分析：}

\vspace{0.2cm}

\noindent
题目先给出了六年级一班女生 $1$ 分钟跳绳的原始成绩，共有 $18$ 个数据：
155，149，116，128，137，151，133，150，145，144，146，123，112，148，139，149，145，125。\\[4pt]
题目要求不是直接计算总数或平均数，而是先按照题干中表格已经给出的五个分段：
100\textasciitilde119、120\textasciitilde129、130\textasciitilde139、140\textasciitilde149、150\textasciitilde159，
把每个成绩分别归入对应区间，再用画“正”字的方法进行整理，最后把每个分段的人数填入统计表。

\vspace{0.4cm}

\noindent\textbf{2. 分类整理：}

\vspace{0.2cm}

\noindent
\textbf{100\textasciitilde119：}116，112，共 \textcolor{blue}{2} 人；\quad 画“正”字可记作：$\vert\vert$。\\
\textbf{120\textasciitilde129：}128，123，125，共 \textcolor{blue}{3} 人；\quad 画“正”字可记作：$\vert\vert\vert$。\\
\textbf{130\textasciitilde139：}137，133，139，共 \textcolor{blue}{3} 人；\quad 画“正”字可记作：$\vert\vert\vert$。\\
\textbf{140\textasciitilde149：}149，145，144，146，148，149，145，共 \textcolor{blue}{7} 人；\quad 画“正”字可记作：正$\vert\vert$。\\
\textbf{150\textasciitilde159：}155，151，150，共 \textcolor{blue}{3} 人；\quad 画“正”字可记作：$\vert\vert\vert$。

\vspace{0.4cm}

\noindent
因此，统计表中“人数”一行应依次填写：\textcolor{blue}{2}，\textcolor{blue}{3}，\textcolor{blue}{3}，\textcolor{blue}{7}，\textcolor{blue}{3}。

\vspace{0.7cm}

\noindent\textbf{3. 绘图（制表）步骤：}

\vspace{0.2cm}

{\small\color{gray}
\textbf{步骤1：}先确定表格结构。表格分为两行六列，左侧第一列为项目栏，其余五列分别对应五个成绩区间。\\
\textbf{步骤2：}在表格上方居中写标题“六年级一班女生1分钟跳绳成绩统计表”，并按照参考答案样式在右上方补写日期“2026年1月”。\\
\textbf{步骤3：}第一行填写“数量/个”和五个成绩分段；第二行左侧填写“人数”。\\
\textbf{步骤4：}把分类统计得到的人数按顺序填入对应单元格：2、3、3、7、3。\\
\textbf{步骤5：}表格线条样式参考前一题，采用上下粗线、中间细线、内部竖分隔线、左右不开口的版式；答案数字使用蓝色，保持与参考答案接近的视觉风格。
}

\vspace{0.8cm}

% ==============================
% 下面开始正式绘制答案表格
% 版式要求：
% 1. 标题居中；
% 2. 日期放在标题下方靠右；
% 3. 表格样式参考本文件前面已有表格：上下粗线、中间细线、内部有竖线；
% 4. 填入的数据使用蓝色，贴近参考答案图片风格。
% ==============================

\begin{center}
{\zihao{4}\bfseries 六年级一班女生 1 分钟跳绳成绩统计表}

\vspace{0.2cm}
\makebox[14.8cm][r]{\textcolor{blue}{2026}\,年\,\textcolor{blue}{1}\,月\hspace*{0.5cm}}

\vspace{0.2cm}
\renewcommand{\arraystretch}{1.6}
\begin{tabular}{c|c|c|c|c|c}
\noalign{\hrule height 1pt}
\makebox[2cm][c]{数量/个} & \makebox[2.2cm][c]{100$\sim$119} & \makebox[2.2cm][c]{120$\sim$129} & \makebox[2.2cm][c]{130$\sim$139} & \makebox[2.2cm][c]{140$\sim$149} & \makebox[2.2cm][c]{150$\sim$159} \\
\hline
人数 & \textcolor{blue}{2} & 3 & \textcolor{blue}{3} & \textcolor{blue}{7} & \textcolor{blue}{3} \\
\noalign{\hrule height 1pt}
\end{tabular}
\end{center}

\vspace{0.6cm}

{\small\color{gray}
\textbf{解答：}\\
先把你原始数据按题目给出的五个区间进行分类，再分别统计每个区间的人数。\\
统计结果是：100\textasciitilde119 有 \textcolor{blue}{2} 人，120\textasciitilde129 有 \textcolor{blue}{3} 人，130\textasciitilde139 有 \textcolor{blue}{3} 人，140\textasciitilde149 有 \textcolor{blue}{7} 人，150\textasciitilde159 有 \textcolor{blue}{3} 人。\\
所以表中应填写为：\textcolor{blue}{2，3，3，7，3}。
}


\newpage



\section{解答：探究总价与长度的正比例关系图像}

\vspace{0.6cm}

% 表格部分
\noindent\textbf{2. 一种花布每米售价 $8$ 元，购买 $2$ 米、$3$ 米$\cdots\cdots$分别需要多少元？}

\vspace{0.3cm}
\begin{center}
\renewcommand{\arraystretch}{1.5}
% 用较粗大一些的表格清晰展示
\begin{tabular}{|c|c|c|c|c|c|c|c|c|c|}
\hline
长度/米 & $1$ & $2$ & $3$ & $4$ & $5$ & $6$ & $7$ & $8$ & $\cdots$ \\
\hline
总价/元 & $8$ 
  & \textcolor{blue!80!black}{\textbf{16}}
  & \textcolor{blue!80!black}{\textbf{24}}
  & \textcolor{blue!80!black}{\textbf{32}}
  & \textcolor{blue!80!black}{\textbf{40}}
  & \textcolor{blue!80!black}{\textbf{48}}
  & \textcolor{blue!80!black}{\textbf{56}}
  & \textcolor{blue!80!black}{\textbf{64}}
  & $\cdots$ \\
\hline
\end{tabular}
\end{center}

\vspace{0.6cm}

% 左右分栏布局：左边题目与解答，右边放 TikZ 正比例函数图
\noindent
\begin{minipage}[t]{0.50\textwidth}
\vspace{0pt} % 贴顶对齐
\normalsize
（1）把上面的表格填写完整。\\
（2）根据表中数据，先在右图中描出长度和总价所对应的点，再把这些点依次连起来。\\
（3）购买这种花布的总价和长度成正比例吗？为什么？

\vspace{0.4cm}
\textcolor{blue!80!black}{\textbf{答：成正比例，因为总价随着长度的变化而变化，且两者的比值一定。}}
\end{minipage}
\hfill
\begin{minipage}[t]{0.45\textwidth}
\vspace{0pt}
\begin{center}
\begin{tikzpicture}[>=Stealth, scale=0.6] % 缩放适配版面右侧

  % 1. 绘制浅蓝色背景网格 (横向9格到x=9，纵向9格到y=9即代表72)
  \draw[step=1.0cm, cyan!50, line width=0.6pt] (0, 0) grid (9.0, 9.0);
  
  % 2. 绘制坐标轴 (黑色带箭头段略长出网格)
  % X 轴
  \draw[->, black, line width=1.0pt] (0, 0) -- (9.8, 0) node[right, font=\small] {长度/米};
  % Y 轴
  \draw[->, black, line width=1.0pt] (0, 0) -- (0, 9.6) node[right, font=\small, xshift=0.1cm] {总价/元};
  
  % 原点刻度 O (0)
  \node[below left, font=\footnotesize] at (0, 0) {$0$};

  % 3. 绘制坐标系横纵轴的刻度文字
  % 横轴：1 到 9
  \foreach \x in {1, 2, ..., 9} {
    \node[below, font=\footnotesize, yshift=-0.1cm] at (\x, 0) {$\x$};
  }
  
  % 纵轴：原题纵轴上对应0-9各格子，1代表8，9代表72
  \foreach \y/\val in {1/8, 2/16, 3/24, 4/32, 5/40, 6/48, 7/56, 8/64, 9/72} {
    \node[left, font=\footnotesize, xshift=-0.1cm] at (0, \y) {\val};
  }

  % 4. 描点连线作图 (答案部分：蓝色折线穿过各个交点)
  \draw[blue!80!black, line width=1.2pt] (1, 1) -- (8, 8);
  
  % 画出实心蓝点即为对应的交点
  \foreach \i in {1, 2, 3, 4, 5, 6, 7, 8} {
    \fill[blue!80!black] (\i, \i) circle (4pt);
  }

\end{tikzpicture}
\end{center}
\end{minipage}

\vspace{0.8cm}
\hrule
\vspace{0.4cm}

{\small\color{gray}
\textbf{解题说明：}\\
1. \textbf{填写表格：}已知花布单价为恒定的 $8\,\text{元/米}$。根据“总价 $=$ 单价 $\times$ 长度”，可以直接依次推算出购买长度 $2, 3, \cdots, 8$ 米所对应的总价分别为 $16, 24, \cdots, 64$ 元。填入表格内标蓝处即可。\\
2. \textbf{描点连线作图：}观察原版题目预设的空白坐标系基座，我们发现其十字网格刻度分配是：横轴单位格代表长度 $1\text{m}$，纵轴单位格恰好代表总价 $8\text{元}$。因此我们将表格上侧列作为横坐标值 $x$，下行对应列金额作为纵轴真实长度换算高度变量 $y$，得到的有序数对 $(1,8)、(2,16) \cdots$ 会极其完美地直接落在背景细网格的各个对角线正交汇处。在此处用明显的实心小圆点深描，再施以直尺连线使其平滑成直线。\\
3. \textbf{正比例函数定性判定解答：}因为在几何图像表面上，这两者的关系被最终绘制成了一条必过原点坐标 $(0,0)$ 的斜向倾斜直线（正比例函数第一象限图像特征）；而在代数根基关系逻辑上：由于总价随着长度成倍递增变化趋势，且商“$\text{总价} \div \text{长度} = \text{单价}(8\text{元/米} > 0)$”，其比率值是一个不可动摇的极常数量，由此彻底证出和满足了判断其正比例属性核心定语：“总价比长度为常数一”。
}

\vspace{0.6cm}

{\small\color{blue!70!black}
\textbf{绘图基本原理与步骤：}\\
\textbf{1. 面向报表排印的矩阵级高分栏：}本题包含了长宽不一的表格以及极其经典的左试卷题干、右描边图系的混合双排版需求。先在顶级流用 \texttt{tabular} 配合 \texttt{\textbackslash textcolor\{blue\}} 高亮醒目突显表格算术推导区，强加阅读辅助反馈。其后排版利用纯正的 \texttt{\textbackslash begin\{minipage\}} 施展极宽版面左右撕裂切割场法：左侧横向切出占空 \texttt{0.50} 承装长文短句试炼题文，右侧空投 \texttt{0.45} 装载包含复杂极坐标矢量的纯粹几何画廊 \texttt{tikzpicture}。它保证了即便后期修改文字，这两片浮动版图也将如同不可相互侵犯的平行宇宙，绝对左右并持且上下顶端齐平约束，呈现出超越回车空格敲击法的几何排版美学。\\
\textbf{2. 自动化步进递推底纹以及降维数轴标记流：}由于原本考题插画底背景有很致密显眼的一大片蔚然蓝海网格，为了防止手绘多条平行线的繁冗，底图一键祭出超级引擎函数 \texttt{\textbackslash draw[step=1.0cm, cyan!50] (0,0) grid (9,9)}，这种暴力网格发生器仅用半行代码就算准出纵横整整十九条长线段相交生成的八十一宫格铺装地基。而在极易写错位的人为数字标注端，放弃手工定位，直接采用自带泛型矩阵性质的 \texttt{\textbackslash foreach} 字典序列扫描。像是那段 \texttt{1/8, 2/16} 构成的映射串列，让纵轴上的标签渲染变成类似解包 Python 列表元组一样，一边抽取出 $y$ 坐标刻度的绝对几何高度 $1,2,3...$，一边向屏幕喷涌抛出渲染出它的实际物理表代表含义 $8,16,24...$ 这不仅缩短了冗长代码，而且防呆容错性使得后续如果物价改为 $10\text{元}$ 也只消在序列库里一键微调全局自动焕发。
}

%% ==============================
%% 第15题：最短路径问题（饮马问题及模型应用）
%% ==============================
\newpage




\section{解答：一架飞机飞行时间和航程}

\vspace{0.4cm}

% 数据表
\begin{center}
\renewcommand{\arraystretch}{1.4}
\begin{tabular}{c|c|c|c|c}
\noalign{\hrule height 1pt}
飞行时间/时 & $2$ & $3$ & $4$ & $6$ \\
\hline
航程/千米 & $1500$ & $2250$ & $3000$ & $4500$ \\
\noalign{\hrule height 1pt}
\end{tabular}
\end{center}

\vspace{0.5cm}

\noindent\textbf{（1）写出几组相对应的航程和飞行时间的比，说一说比值表示什么。}

\vspace{0.3cm}

\noindent
$\dfrac{1500}{2} = 750$，\quad
$\dfrac{2250}{3} = 750$，\quad
$\dfrac{3000}{4} = 750$，\quad
$\dfrac{4500}{6} = 750$

\vspace{0.3cm}

\noindent
比值都是 $750$，表示飞机的\textbf{\textcolor{red!70!black}{速度}}，即每小时飞行 $750$ 千米。

\vspace{0.6cm}

\noindent\textbf{（2）航程和飞行时间成正比例吗？为什么？}

\vspace{0.3cm}

\noindent
\textbf{答：}成正比例。

\vspace{0.2cm}

\noindent
因为：航程 $\div$ 飞行时间 $= 750$（定值），即两个量的\textbf{比值一定}，\\[4pt]
所以航程和飞行时间成\textbf{\textcolor{red!70!black}{正比例}}。

\vspace{0.6cm}

\noindent\textbf{（3）在下图中描出飞行时间和航程所对应的点，再顺次连接起来。}

\vspace{0.4cm}

\begin{center}
\begin{tikzpicture}[>=Stealth, x=1cm, y=0.5cm] % 整体思路：先建立“时间-航程”坐标系，再画网格和刻度，随后把题目数据按比例落点，最后连成折线/直线，直观看出正比例图像过原点且各点共线
  % 网格（长方形网格，长宽比 2:1）
  % x方向每1单位 = 1cm，y方向每1单位 = 0.5cm
  \draw[step=1, xstep=1, ystep=1, gray!40, thin] (0,0) grid (8,6); % draw 表示绘制图形；step=1 设总步长为 1；xstep=1 与 ystep=1 分别指定横纵网格间距都按坐标单位 1 生成；gray!40 表示 40% 深度的灰色；thin 表示细线；(0,0) grid (8,6) 表示从原点到右上角 (8,6) 生成整片网格

  % 坐标轴
  \draw[->, line width=1pt] (0,0) -- (8.5,0) node[right] {\small 飞行时间/时}; % -> 表示终点带箭头；line width=1pt 设置坐标轴粗细；(0,0)--(8.5,0) 画 x 轴；node[right] 在终点右侧放标签；\small 控制标签字号
  \draw[->, line width=1pt] (0,0) -- (0,7.5) node[above] {\small 航程/千米}; % 同理画 y 轴；终点设在 (0,7.5) 留出轴标题位置；node[above] 让标题放在线段终点上方

  % x 轴刻度标注
  \foreach \x in {0,1,...,8} { % foreach 循环变量 \x 依次取 0 到 8，用于批量画横轴刻度与数字
    \draw[line width=0.6pt] (\x, -0.15) -- (\x, 0); % 在每个整数横坐标处画一条短刻度线；line width=0.6pt 表示刻度线较细；从 (\x,-0.15) 到 (\x,0) 是一条竖向短线
    \node[below, font=\small] at (\x, -0.15) {$\x$}; % 在刻度下方放数字；below 表示节点在参考点下方；font=\small 指定字号；at (\x,-0.15) 指定位置；{$\x$} 显示当前循环值
  } % 结束 x 轴刻度循环

  % y 轴刻度标注（每格 = 750千米）
  \foreach \y/\lab in {1/750, 2/1500, 3/2250, 4/3000, 5/3750, 6/4500} { % 这里采用“坐标值/显示标签”成对循环，\y 是纵坐标，\lab 是要写出的实际航程数
    \draw[line width=0.6pt] (-0.1, \y) -- (0, \y); % 在 y 轴左侧画短刻度；x 从 -0.1 到 0，形成横向小短线
    \node[left, font=\small] at (-0.1, \y) {$\lab$}; % 在刻度左边标上对应航程；left 表示文字出现在点左侧；\lab 显示 750、1500 等实际数值
  } % 结束 y 轴刻度循环

  % 从原点出发的线段，顺次连接所有数据点
  % 数据点：(0,0), (2,2), (3,3), (4,4), (6,6)
  \draw[red!70!black, line width=1.5pt] (0,0) -- (2,2) -- (3,3) -- (4,4) -- (6,6); % 用红色较粗折线依次连接原点和各数据点；颜色 red!70!black 更醒目；line width=1.5pt 强调关系线；连续 -- 表示折线逐点相连

  % 标注各点数值
  \node[above left, font=\scriptsize, red!70!black] at (2,2) {$(2,\,1500)$}; % 在点 (2,2) 左上方写坐标说明；above left 指节点位于左上；font=\scriptsize 用更小字号避免拥挤；\, 插入细空格分隔两个数
  \node[above left, font=\scriptsize, red!70!black] at (3,3) {$(3,\,2250)$}; % 在第二个数据点附近标注对应“时间、航程”数对
  \node[above left, font=\scriptsize, red!70!black] at (4,4) {$(4,\,3000)$}; % 在第三个数据点附近标注数据
  \node[above left, font=\scriptsize, red!70!black] at (6,6) {$(6,\,4500)$}; % 在第四个数据点附近标注数据

  % 各数据点的实心圆点
  \foreach \p in {(2,2),(3,3),(4,4),(6,6)} { % 用循环统一处理四个数据点，\p 每次代表一个坐标对
    \fill[red!70!black] \p circle (2.5pt); % fill 表示填充图形；颜色同折线；\p circle (2.5pt) 表示以当前点为圆心画半径 2.5pt 的实心圆点
  } % 结束数据点填充循环

\end{tikzpicture} % 结束“飞机飞行时间和航程”坐标图
\end{center}

\vspace{0.4cm}

{\small\color{gray}
\textbf{说明：}所有点在同一条过原点的直线上，这正是正比例关系的图像特征。}

\vspace{0.6cm}

{\small\color{blue!70!black}
\textbf{绘图基本原理与步骤：}\\
\textbf{1. 坐标系搭建}：通过 \texttt{x=1cm, y=0.5cm} 非等比缩放建立矩形网格坐标系。X 轴表示飞行时间（时），Y 轴表示航程（千米）。使用 \texttt{->} 箭头语法绘制带标签的坐标轴。\\
\textbf{2. 刻度映射}：Y 轴每格代表 750 千米，通过 \texttt{foreach} 循环生成刻度标注。数据点坐标按 $y = \text{航程} \div 750$ 映射（如 $1500 \to y=2$，$4500 \to y=6$）。\\
\textbf{3. 数据可视化}：使用 \texttt{draw[red!70!black]} 连接所有数据点成折线，\texttt{fill} 绘制实心圆点，\texttt{node[above left]} 附挂坐标标签。所有点共线过原点，直观展示正比例关系。
}


%% ==============================
%% 连线题：数字中“5”的意义
%% ==============================
\newpage
